
\documentclass[11pt]{article}

\usepackage{calc}
\usepackage{graphics}
%\usepackage{latex8}
\usepackage{times}
\usepackage{epsf}
\usepackage{epsfig}
\usepackage{graphicx}
\usepackage{changebar}
\usepackage{portland}
\usepackage{lscape}

\setlength{\textheight}{8.50in}
\setlength{\textwidth}{6.5in}
\setlength{\topmargin}{-0.3in}
%\setlength{\leftmargin}{2.9in}
%\setlength{\rightmargin}{-2.9in}
\setlength{\oddsidemargin}{0in}
\setlength{\parindent}{0.5in}
\setlength\parindent{0in}
\setlength\parskip{0.1in} 

\newcommand{\note}[1]{{$\rightarrow$ \bf Note: \emph{#1}}}

\begin{document}

\title{
\includegraphics[height=1.5in]{v3vee.pdf}
\includegraphics[height=1.5in]{logo6.png} \\
\vspace{0.5in} 
Palacios Internal Developer Manual
}


\maketitle


This manual is written for internal Palacios developers. It contains
information on how to obtain the Palacios code base, how to go about
the development process, and how to commit those changes to the
mainline source tree.  This assumes that the reader has read {\em An
Introduction to the Palacios Virtual Machine Monitor -- Release 1.0}
and also has a slight working knowledge of {\em git}.  You will also
want to read the document {\em Building a bootable guest image for
Palacios and Kitten} in order to understand how to build an extremely
lightweight guest environment, suitable for testing.

Please note that Palacios and Kitten are under rapid development,
hence this manual may very well be out of date!

\newpage
\tableofcontents
\newpage
\listoffigures
\newpage

\section{Overview}


Both Palacios and Kitten follow a hybrid development process that
uses both the centralized repository and distributed development
models. A central repository exists that holds the master version of
the code base. This central repository is cloned by multiple people
and in multiple places to support various development efforts. A
feature of \texttt{git} is that every developer actually has a fully copy of
the entire repository, and so can function independently until such
time as they need to re-sync with the master version. 

There are typically multiple levels of access to the central
repository, that are granted based on the type of developer being
granted access. The three basic developer types and their access
privileges are:

\begin{itemize}
\item Core developers: These are the lead developers and are in
charge of managing the master repository. They have full read/write
access permissions to the central repository.

\item Internal developers: Formal members of the development
team. These people are capable of pulling directly from the central
repository, but lack the ability to write directly to it. 

\item External developers: People who are not actual members of the
development team. These people can only access the public repository
which is only updated to contain the release versions. 
\end{itemize}

Students doing independent study or REUs related to Palacios are set
up as internal developers.  EECS 441 (Resource Virtualization)
students are generally either set up as internal or external
developers, depending on their projects.

Because internal and external developers cannot write directly to the
master repository, they need to first submit their changes to a core
developer before they can be added to the mainline. We will discuss
that process in Section~\ref{sec:submission}.


\section{Checking out Palacios}

The central Palacios repository is located on {\em
newskysaw.cs.northwestern.edu} in {\em /home/palacios/palacios}. All
internal developers have read access to the directory. Each developer
must create their own local version of the repository, this is done
with {\em git clone}.

\begin{verbatim}
git clone /home/palacios/palacios
\end{verbatim}

On the machine {\em newbehemoth.cs.northwestern.edu} you will want to
use the following command instead. The {\em newskysaw} home
directories are NFS-mounted on /home-remote.


\begin{verbatim}
git clone /home-remote/palacios/palacios
\end{verbatim}


On any other machine, you can clone the repository via ssh, provided
you have a newskysaw account:
 
\begin{verbatim}
git clone ssh://you@newskysaw.cs.northwestern.edu//home/palaicos/palacios
\end{verbatim}

External developers can clone the public repository via the web.  The
public repository tracks the release and main development branch
(e.g., devel) of the internal repository with a 30 minute delay.  The
access information is available on the web site (http://v3vee.org).
The web site also includes a publically accessible gitweb interface to
allow browsing of the public repository.

This creates a local copy of the repository at {\em ./palacios/}.

Note that both {\em newskysaw} and {\em newbehemoth} have all the
tools installed that are needed to build and test Palacios and Kitten.
If you develop on another machine, you will need to set those tools up
for yourself.  This isn't hard and the tools are all free.  See the
technical report for what tools you will need.

When you first clone the repository, you will get the {\em master}
branch, which is used to generate releases.   All development work is
done in the {\em devel} branch of the repository. The developer can
access this branch via:

\begin{verbatim}
git checkout --track -b devel origin/devel
\end{verbatim}

or 

\begin{verbatim}
/opt/vmm-tools/bin/checkout_branch devel
\end{verbatim}

{\em Important:}
Note that Palacios is very actively developed so the contents of the
{\em devel} branch are frequently changing. In order to keep up to
date with the latest version, it is necessary to periodically pull the
latest changes from the master repository by running \verb.git pull..



\section{Checking out Kitten}

Kitten is available from Sandia National Labs, and is the main host OS
we are targeting with Palacios. Loosely speaking, core Palacios
developers are internal Kitten developers, and internal Palacios
developers are external Kitten developers. The public repository for
Kitten is at {\em http://code.google.com/kitten}.  To simplify things,
we are maintaining a local mirror copy in {\em /home/palacios/kitten}
that tracks the public repository.

Kitten uses Mercurial for their source management, so you will have to
make sure the local mercurial version is configured correctly.
Specifically you should add the following Python path to your shell environment.

\begin{verbatim}
export PYTHONPATH=/usr/local/lib64/python2.4/site-packages/
\end{verbatim}

You can then clone Kitten from the local mirror.   On {\em newskysaw},
run: 
\begin{verbatim}
hg clone /home/palacios/kitten
\end{verbatim}
On {\em newbehemoth}, run
\begin{verbatim}
hg clone /home-remote/palacios/kitten
\end{verbatim}
On other machines, run
\begin{verbatim}
hg clone ssh://you@newskysaw.cs.northwestern.edu//home/palacios/kitten
\end{verbatim}


Both the Kitten and Palacios clone commands should be run from the
same directory. This means that both repositories should be located at
the same directory level. The Kitten build process depends on this.

{\em Important:} Like Palacios, Kitten is under active development,
and its source tree is frequently changing. In order to keep up to
date with the latest version, it is necessary to periodically pull the
latest changes from the mirror repository by running \verb.hg.
pull. followed by \verb.hg update..

\section{Compiling Palacios}

The Palacios build process has been changed recently from a homegrown
environment to the widely used KBuild environment.  KBuild is also
used for building Kitten.  Because KBuild is the build environment
used for Linux, much of what you learn about configuring and building
Linux kernels is readily applicable to Palacios and Kitten. 

The output of the Palacios build process is a static library that
includes the Palacios VMM and relevant guest support code blocks. This
static library is then linked into a host operating system. Palacios
internally supports GeekOS and can generate a complete OS image via a
unified build process.  By complete OS image, we mean an ISO image
containing GeekOS, Palacios, and a guest image (another ISO image) for
testing. 

These days, however, Palacios is typically embedded into Kitten, not
GeekOS.  To combine Palacios with Kitten, it is necessary to first
configure and build Palacios, then to configure and build Kitten,
linking in Palacios.   Kitten can also be configured to link in a
guest image for testing.  

\subsection{Configuration}

To configure Palacios, enter the top level Palacios directory and
execute:
\begin{verbatim}
make clean
make xconfig
\end{verbatim}
At this point, you will see be presented with a KBuild configuration
screen, similar to what you would see in configuring a Linux kernel.
Palacios has far fewer options, however.   If you don't have X or
don't want the graphical configuration system, you can also use the
\verb.menuconfig. or \verb.config. targets.  The available options
change over time, so we do not cover all of them here, but here are a
few that are usually important, with their recommended values noted:
\begin{itemize}
\item Target Configuration:   
\begin{itemize}

\item Red Storm (Cray XT3/XT4) --- turn on to target Cray XT4
supercomputers.  (off)
\item AMD SVM Support --- targets AMD processors with the SVM hardware
virtualization features (on)
\item Intel VMX support --- targets Intel processors with the VMX
hardware virtualization features (on)
\item Compile for a multi-threaded OS (on)
\item Enable VMM telemetry support --- this is lightweight logging and
data collection (on)  
\item Enable VMM instrumentation --- this is heavyweight logging and
data collection (off)
\item Enable passthrough video --- this lets a guest write directly to
the video card (on)
\item Enable experimental options --- this makes it possible to select
features that are under current development (on).  You probably want
to leave VNET turned off.  VNET is an experimental VMM-embedded
overlay network under development by Lei Xia and Yuan Tang.
\item Enable built-in versions of stdlib functions --- this adds
needed stdlib functions that the host OS may not supply.  For use with
Kitten turn on and enable strncasecmp() and atoi().
\item Enable built-in versions of stdio functions (off)
\end{itemize}
\item Symbiotic Functions (these are experimental options for Jack
Lange's thesis). 
\begin{itemize}
\item Enable Symbiotic Functionality --- This adds symbiotic features
to Palacios, specifically support for discovery and configuration by
symbiotic guests, the SymSpy passive information interface for
asynchronous symtiotic guest $\leftrightarrow$ symbiotic VMM
information flow, and the SymCall functional interface for synchronous
symbiotic VMM $\rightarrow$ symbiotic guest upcalls.  (on)
\item Symbiotic Swap --- Enables the SwapBypass symbiotic service for
symbiotic Linux guests.  (off)
\end{itemize}
\item Debug Configuration
\begin{itemize}
\item Compile with Debug info --- adds debug symbols (-g)  (off)
\item Enable Debugging --- makes it possible to show PrintDebug output
(on).   You can selectively turn on debugging output for each major
VMM component, including shadow paging, nested paging, control
registers, interrupts, I/O, instruction emulation and XED, halt, and
the device manager.   Note that the more debugging output you turn on,
the slower the VMM will go since it will have to wait for the prints
to finish. 
\end{itemize}
\item BIOS Selection --- Lets you select which code blobs will be
used for bootstrapping the guest.   There are currently three:  a
BIOS, a Video BIOS, and the VMXAssist V8086 service (the latter is used only on
Intel VMX).   Generally, you should not need to change these.
\item Virtual Devices --- virtual devices can be instantiated and
added to a guest.   The following is a list of the currently
implemented virtual devices.
\begin{itemize}
\item BOCHS Debug Console Device --- used for debugging output from
the guest BIOS (on)
\item OS Debug Console Device --- used for debugging output from
the guest kernel (on)
\item 8259A PIC - legacy Programmable Interrupt Controller chip --- used
for bootstrap of most guests (on)
\item APIC - in-processor Advanced Programmable Interrupt Controller --- used
for interrupt delivery on almost all guests (on)
\item IOAPIC - Off-chip APIC  --- used for interrupt deliver for almost
all guests (on)
\item i440fx Northbridge --- emulation of a typical PC North Bridge
chip, used on almost all guests (on)
\item PIIX3 Southbridge --- emulation of a typical PC South Bridge
chip, used on almost all guests (on)
\item PCI --- emulation of a PCI bus - needed for attaching most
devices (on)   
\begin{itemize}
\item Passthrough PCI --- allows us to make a hardware PCI device visible and
directly accessible by the guest (on)
\end{itemize}
\item Generic --- this is a run-time configurable device that can intercept I/O port read/writes and memory region reads/writes.   Intercepted reads and writes can either be ignored or forwarded to actual hardware, and the data flow can optionally be printed.   This is a useful tool with at least three purposes.  First, it makes it possible to ``stub out'' hardware that isn't currently implemented and for which we don't want to allow passthrough access. Second, it makes it possible to provide low-level passthrough access to physical hardware.   Third, it can be used to spy on guest/device interactions, which is very helpful when trying to understand the interface of a device.
\item NVRAM - motherboard configuration memory --- needed by BIOS bootstrap (on)
\item Keyboard - Generic PS/2 keyboard, including mostly broken mouse
implementation (on)
\item IDE --- Support for virtual IDE controllers that support disks
and CD ROMs (on)

\item NE2K - NE2000 and RTL8139 network devices (off)
\item CGA - CGA video card (paritial implementation) (off)
\begin{itemize}
\item Telnet Virtual Console (off)   When CGA and Telnet Console are
on, it is possible to telnet to the console of the guest.   Eventually
the rest of this will provide simple bitmapped video console for VNC
access.
\end{itemize}
\item RAMDISK storage backend --- used to create RAM disk
implementations of block devices (on)
\item NETDISK storage backend --- used to create network-attached disk
implementations of block devices, e.g., network block devices (on)
\item TMPDISK storage backend --- used to create temporary storage
implementations of block devices (on)
\item Linux Virtio Balloon Device --- used for memory ballooning by
Linux virtio-compatible guests (on)
\item Linux Virtio Block Device --- used for fast block device support
by Linux virtio-compatible guests (on)
\item Linux Virtio Network Device --- used for fast network device support
by Linux virtio-compatible guests (on) 
\item Symbiotic Swap Disk (multiple versions) --- used for the
SwapBypass service (off)
\item Disk Performance Model --- used for the
SwapBypass service (off)
\end{itemize}
\end{itemize}

\subsection{Compilation}

After configuring Palacios---remember to save your changes---you can
compile it by executing
\begin{verbatim}
make 
\end{verbatim}
This will produce the file {\em libv3vee.a} in the current directory.
This static library contains the Palacios VMM and is ready for
embedding into an OS, such as Kitten.  The library provides the
ability to instantiate and run virtual machines.  By default, on a 64
bit machine, the library is compiled for 64 bit machines (x86\_64),
while on a 32 bit machine, it is compiled for 32 bit machines.  You
can override this using the ARCH=i386 or ARCH=x86\_64 arguments to the
make, provided you have the relevant tools available.  The 64 bit
version is what you need for use with Kitten.  A 64 bit Palacios can
run both 64 and 32 bit guests.  Both {\em newskysaw} and {\em
newbehemoth} are 64 bit machines. 


\section{Compiling Kitten}
Kitten requires a 64-bit version of Palacios, so make sure that
Palacios has been correctly compiled before compiling Kitten.  The
current default for Palacios is 64 bit. 

\subsection{Configuration}
Kitten borrows a lot of concepts from Linux, including the Linux build
process. As such it must be configured before it is actually compiled.
The Kitten configuration process is the same as Linux, and can be
accessed via any of these make targets.
\begin{itemize}
\item \verb.make xconfig.
\item \verb.make config.
\item \verb.make menuconfig.
\end{itemize}

Of course, there are a range of configuration options.  In the
following, we note only the most important:
\begin{itemize}
\item Target Configuration
\begin{itemize}
\item System Architecture --- you probably want to set this to
PC-Compatible, unless you are working on Red Storm.
\item Processor Family --- you want to set this to either
AMD-Opteron/Athlon64 or Intel-64/Core2, depending on whether you have
a 64 bit AMD or 64 bit Intel processor.
\end{itemize}
\item Virtualization
\begin{itemize}
\item Include Palacios VMM --- this will link against the Palacios
library (on)
\item Path to pre-built Palacios tree --- directory where libv3vee.a
can be found.
\item Path to guest image --- location of the test guest OS
image that will be embedded.  We will say more about this later.
Essentially, however, a guest image consists of a blob that begins
with an XML description of the desired guest environment and the
contents of the remainder of the blob.   The remainder of the blob
usually contains disk or cd images.
\end{itemize}
\item Networking 
\begin{itemize}
\item Enable LWIP TCP/IP stack.  This activates a simple TCP/IP stack
that things like NETDISK can use. (on)
\end{itemize}
\item Device Drivers
\begin{itemize}
\item VGA Console --- driver for basic video.  If you turn on
passthrough video in Palacios, you should turn this off.
\item Serial Console --- driver for serial port console.  (on)
\item VM Console --- driver for Kitten console on top of Palacios.  If
Kitten is run {\em as a guest}, and it has VM Console on, then it can output
cleanly via the Palacios OS Console device (off).
\item NE2K Device Driver --- driver for NE2K and RTL8139 network cards
(on)
\item VM Network Driver --- driver for Kitten network output using
Palacios.  If Kitten is run {\em as a guest}, and it has VM Network
Driver, then it can send and receive packets using the Palacios Linux
virtio network device.  (off)
\end{itemize}
\item ISOIMAGE configuration:  Kitten kernel arguments.   Note that
this is NOT for the guest image, but rather for the Kitten image.  You
can leave this alone.  For Palacios operation, it's important that the
option \verb.console=serial. appears.  If the NE2K/RTL8139 driver
should be used \verb.net=rtl8139. should appear. 
\item Kernel Hacking
\begin{itemize}
\item Kernel Debugging --- here you can turn on various Kitten
Linux-like debugging features.   Only a few are noted below:
\begin{itemize}
\item Compile the kernel with debug info --- if this is on, you will
have debugging information compiled in (-g)
\item KGDB --- if you have this enabled, you will be able to attach to
the running kernel from the GDB debugger.  This means you can also
attach to Palacios, which is embedded.   If you want to debug Palacios
using KGDB, be sure to turn on debugging in Palacios as well.
\end{itemize}
\end{itemize}
\item Include Linux compatability layer --- if this is on, you can 
selectively add Linux system calls and other functionality to Kitten.
Kitten is able to run Linux ELF executables as user processes with
this layer.   
\end{itemize}

The guest OS that is to be booted as a VM is included as a blob
pointed to by ``Path to guest image''.   The blob starts with an XML
description of the guest, followed by other chunks of data used, for
example, as the content of virtual hard drives or CD ROMs.  Please see
Section~\ref{sec:guestconfig} for basic information on how to use the
guest builder to assemble a guest OS blob. 

By default, the init task that is executed after Kitten boots (located
in user/hello\_world) does a number of Kitten tests.  One of these is
a test of the VMM API, which is implemented using Palacios.  When this
test is done, a VM is created, configured according to the XML, and
the guest OS blob is launched in it.


\subsection{Compilation}

After Kitten has been configured it can be compiled.  This is done
simply by executing
\begin{verbatim}
make isoimage
\end{verbatim}
This command will compile Kitten (with Palacios embedded in it) and
the init task (which will contain the guest OS blob), and then
assemble an ISO image file which can used to boot a machine.  The ISO
image is located at {\em ./arch/x86\_64/boot/image.iso}.  

This image file can be used for booting a QEMU emulation environment,
for booting a remote machine using PXE, or can be burned to CD/DVD for
booting a machine physically. 


\section{Basic Guest Configuration}
\label{sec:guestconfig}

To configure a guest, you write an XML configuration file, which
contains references to other files that contain data needed to
instantiate stateful devices such as virtual hard drives and CD ROMs.
You supply this information to a guest builder utility that assembles
a guest image suitable for reference in the Kitten configuration, as
described above.  

The guest builder utility is located in {\em
palacios/utils/guest\_creator}.  You will need to run \verb.make. in that
directory to compile it, resulting in the executable named {\em
build\_vm}.  Also located in that directory is an example configuration
file, named {\em default.xml}.   We typically use this file as a
template.  It is carefully commented.  In summary, a configuration
consists of
\begin{itemize}
\item Physical memory size of the gust
\item Basic VMM settings, such as what form of virtual paging is to be
used, the scheduler rate, whether services like telemetry are on, etc.
\item A memory map that maps regions of the host physical address
space to the guest physical address space.  This can, for example,
make a framebuffer visible in the guest.
\item A list of the files that will be used in assembling the image.
For example, the contents of a boot CD.
\item A list of the devices that the guest will have, including
configuration data for each device.
\end{itemize}
There are a few subtlies involved with devices.  One is that some
devices form attachement points for other devices.  The PCI device is
an example of this.  Another is that each device needs to define how
it is attached (e.g. direct (implicit), via a PCI bus, etc.)
Finally, there may be multiple instances of devices.   For example, a
PCI passthrough device is instantiated for every underlying PCI device
we want to make visible in the guest. 

The XML configuration format is carefully designed to be extensible.
For example, new devices could use additional or new configuration
options.  The configuration parser in Palacios essentially ignores XML
blocks it doesn't understand. 


To build a guest, one runs
\begin{verbatim}
palacios/utils/guest_creator/build_vm myconfig.xml -o myimage.dat
\end{verbatim}
Here, {\em myimage.dat} is the guest image that can be given to
Kitten. 

A common kind of guest used for testing is one that boots some form of
bootable Linux distribution, or other a live OS distribution.  These
distributions are CD ROM images (ISOs).  A range of them are available
on {\em newskysaw} under {\em /opt/vmm-tools/isos}.  We often use
Puppy Linux ({\em puppy.iso}) or Finnix ({\em finnix.iso}), for
example, but isos are also available for Windows of different flavors,
DOS, GeekOS, and others.  If you just want to use some guest ISO image
like this, you can generally just copy the default XML file, and
modify the
\verb.filename=. attribute here:
\begin{verbatim}
 <files>
    <!-- The file 'id' is used as a reference for 
         other configuration components -->
    <file id="boot-cd" filename="/home/jarusl/image.iso" />
    <!--<file id="harddisk" filename="firefox.img" />-->
 </files>
\end{verbatim}

For careful, repeatable experimentation, it is often convenient to
build your own simplified Linux guest image.  It will boot {\em much}
faster than a full blown distribution and you can readily set up an
environment in which you can exert very tight control, being able to
modify the Linux kernel, the included files (e.g., benchmarks), and
other components very rapidly.  To learn more about how to do this,
please consult the separate document named {\em Building a bootable
guest image for Palacios and Kitten}.

\section{Running Palacios/Kitten}
Kitten and Palacios are capable of running under QEMU, which makes
debugging much simpler.  QEMU is a user-level Linux or Windows program
that emulates a PC machine. 

The basic form of the command to start the QEMU emulator is:
\begin{verbatim}
/usr/local/qemu/bin/qemu-system-x86_64 -smp 1 -m 1024 \
        -serial file:./serial.out \
        -cdrom ./arch/x86_64/boot/image.iso  \
        < /dev/null
\end{verbatim}

The command starts up a single processor emulated machine, with 1GB of
RAM and a CD-ROM drive loaded with the Kitten ISO image.  All output
to the serial port is written directly to a file called {\em
  serial.out}. This command can be copied into a shell script for easy
access.

We can also run Palacios/Kitten on physical hardware.  The slow way is
to burn the Kitten ISO image onto a CD ROM and then boot the test
machine with it.  The much faster way is to set the test machine up to
use the PXE network boot system (most modern BIOSes support this), and
boot your Kitten image over the network.  The debugging output will
then appear on the actual serial port of the physical machine.  For
the Northwestern environment, please talk to Jack Lange or Peter Dinda
if you need to be able to do this.  Northwestern has a range of AMD
and Intel boxes for testing, as do UNM and Sandia.    A different form
of network boot is used for Red Storm. 


\section{Development Guidelines}

There are standard requirements we have for code entering the mainline. 

First and foremost, Palacios is designed to be OS independent and
support 32-bit and 64-bit architectures. This means that developers should
not include any external OS specific dependencies in any Palacios
component. Also all changes need to be tested on both 32-bit and 64-bit
architectures to make sure that they compile as well as run correctly.

\paragraph*{Coding Style}

``The use of equal negative space, as a balance to positive space, in a
composition is considered by many as good design. This basic and often
overlooked principle of design gives the eye a 'place to rest,'
increasing the appeal of a composition through subtle means.''
\newline\newline
Translation: Use the space bar, newlines, and parentheses. 

Curly-brackets are not optional, even for single line conditionals. 

Tabs should be 4 characters in width.

{\em Special:} If you are using XEmacs add the following to your \verb!init.el! file:
\begin{verbatim}
(setq c-basic-offset 4)
(c-set-offset 'case-label 4)
\end{verbatim}

{\em Bad}
\begin{verbatim}
if(a&&b==5||c!=0) return;
\end{verbatim}


{\em Good}
\begin{verbatim}
if (((a) && (b == 5)) || 
    (c != 0)) {
       	return;
}
\end{verbatim}



\paragraph*{Fail Stop}
Because booting a basic Linux kernel results in over 1 million VM exits
catching silent errors is next to impossible. For this reason
ANY time your code has an error it should return -1, and expect the
execution to halt. 

This includes unimplemented features and unhandled cases. These cases
should ALWAYS return -1. 


\paragraph*{Function names}
Externally visible function names should be used rarely and have
unique names. Currently we have several techniques for achieving this:

\begin{enumerate}
\item \verb.#ifdefs. in the header file
\newline
When the V3 Hypervisor is compiled it defines the symbol
\verb.__V3VEE__. Any function that is not needed outside the Hypervisor
context should be inside an \verb.#ifdef __V3VEE__. block, this will make it
invisible to the host environment.

\item Static Functions
\newline
Any utility functions that are only needed in the .c file where they
are defined should be declared as static and not included in the
header file. You should make an effort to use static functions
whenever possible. 

\item \verb.v3_. prefix \newline Major interface functions should be
  named with the prefix \verb.v3_. This allows easy understanding of
  how to interact with the subsystems.  In the case that they need to
  be externally visible to the host OS, make them unlikely to collide
  with other functions.
\end{enumerate}

\paragraph*{Debugging Output}
Debugging output is sent through the host OS via functions in the
\verb.os_hooks. structure. These functions have various wrappers of the form
\verb.Print*., with \texttt{printf}-style semantics. 

Two functions of note are \verb.PrintDebug. and \verb.PrintError..

\begin{itemize}

\item PrintDebug:
\newline
Should be used for debugging output that will often be turned off
selectively by the VMM configuration.

\item PrintError
\newline
Should be used when an error occurs, this will never be optimized out
and will always print. 
\end{itemize}


\section{Code Submission}
\label{sec:submission}

To commit changes to the central repository they need to be exported
as a patch set that can be applied directly to a mainline. Both Git
and Mercurial contain functionality to allow developers to maintain
changes as a patch set. There are also a few options that make dealing
with patches easier.

\subsection{Palacios}

Git includes support for directly exporting local repository commits
as a patch set. The basic operation is for a developer to commit a
change to a local repository, and then export that change as a patch
that can be applied to another git repository. While this is
functionally possible, there are a number of issues. The main problem
is that it is difficult to fully encapsulate a new feature in a single
commit, and dealing with multiple patches that often overwrite each
other is not a viable option either. Furthermore, once a patch is
applied to the mainline, it will generate a conflicting commit that
will become present when the developer next pulls from the central
repository. This can result in both repositories getting out of
sync. It is possible to deal with this by manually re-basing the local
repository, but it is difficult and error-prone. 

This approach also does not map well when patches are being revised. A
normal patch will go through multiple revisions as it is reviewed and
modified by others. This often leads to synchronization issues as well
as errors with patch revisions. Ultimately it is the responsibility of
the developer to generate a patch that will apply cleanly to the
mainline.

For this reason most internal developers should seriously consider
{\em stacked git}. Stacked git is designed to make patch development
easier and less of a headache. The basic mode of operation is for a
developer to initialize a patch for a new feature and then continuously
apply changes to the patch. Stacked Git allows a developer to layer a
series of patches on top of a local git repository, without causing
the repository to unsync due to local commits. Basically, the
developer never commits changes to the repository itself but instead
commits the changes to a specific patch. The local patches are managed
using stack operations (push/pop) which allows a developer to apply
and unapply patches as needed. Stacked git also manages new changes to
the underlying git repository as a result of a pull operation and
prevents collisions as changes are propagated upstream. For instance
if you have a local patch that is applied to the mainline as a commit,
when the commit is pulled down the patch becomes empty because it is
effectively identical to the mainline. It also makes incorporating
external revisions to a patch easier. Stacked git is installed on {\em
newskysaw} in \verb./opt/vmm-tools/bin/. 

Brief command overview:
\begin{itemize}
\item \verb.stg init. -- Initialize stacked git in a given branch
\item \verb.stg new. -- create a new patch set, an editor will open
asking for a commit message that will be used when the patch is
ultimately committed.
\item \verb.stg pop. -- pops a patch off of the source tree.
\item \verb.stg push. -- pushes a patch back on to a source tree.
\item \verb.stg export. -- exports a patch to a directory as a file
that can then be emailed.
\item \verb.stg refresh. -- commits local changes to the patch set at
the top of the applied stack.
\item \verb.stg fold. -- Apply a patch file to the current
pat1ch. (This is how you can manage revisions that are made by other developers).
\end{itemize}

You should definitely look at the online documentation to better
understand how stacked git works. It is not required of course, but if
you want your changes to be applied its up to you to generate a patch
that is acceptable to a core developer. Ultimately using Stacked git
should be easier than going it alone.


All patches should be emailed to Jack for inclusion in the
mainline. An overview of the organization is given in
Figure~\ref{fig:process}. You should assume that the first revision of
a patch will not be accepted, and that you will have to make
changes. Furthermore, the final form of the patch most likely will not
be exactly what you submitted. 

 
\begin{figure}[t]
\begin{center}
\includegraphics[height=3.5in]{dev_chart.pdf}
\end{center}
\caption{Development organization}
\label{fig:process}
\end{figure}


\subsection{Kitten}

Writing code for Kitten follows essentially the same process as
Palacios. The difference is that the patches need to be emailed to the
Kitten developers. To send in a patch, you can just email it to the
V3Vee development list.


Also, instead of Stacked git you should use Mercurial patch
queues. This feature is enabled in your .hgrc file.
\begin{verbatim}
[extensions]
hgext.mq=
\end{verbatim}

Mercurial queues use the same stack operations as stacked git, however
it does not automatically handle the synchronization with pull
operations. Before you update from the central version of Kitten you
need to pop all of the patches, and then push them once the update is
complete.

Basically:
\begin{verbatim}
hg qpop -a
hg pull
hg update
hg qpush -a
\end{verbatim}


%Also, remember that Kitten is not a Northwestern project and is being
%developed by professional developers at Sandia National Labs. So keep
%in mind that you are representing Northwestern and the rest of the
%Palacios development group. We are collaborating with them because
%Kitten and the resources they have are very important for our research
%efforts. They are collaborating with us because they believe that
%Palacios might be able to help them. Therefore it is important that we
%continue to ensure that they see value in our collaboration. In plain
%terms, we need to make sure they think we're smart and know what we're
%doing. So please keep that in mind when dealing with the Kitten group.


\section{Networking}

Both the Kitten and GeekOS substrates on which Palacios can run
currently include drivers for two simple network cards, the NE2000,
and the RTL8139.  The Kitten substrate is acquiring an ever increasing
set of drivers for specialized network systems.   A lightweight
networking stack is included so that TCP/IP networking is possible
from within the host OS kernel and in Palacios.  

When debugging Palacios on QEMU, it is very convenient to add an
RTL8139 card to your QEMU configuration, and then drive it from within
Palacios.  QEMU can be configured to provide local connectivity to the
QEMU emulated machine, including bridging the emulated machine with a
physical network.  Local connectivity can be done with redirection, or
with a TAP interface.  For global connectivity, a TAP interface must
be used; it is bridged to a physical interface.

\section{Configuring the development host's QEMU network}

To get local connectivity with redirection, no networking changes on
the host are needed.  However, people usually want to use TAP-based
networking, which does require changes.  For one thing, TAP interfaces
can be inspected with tools like wireshark, which makes for much
easier debugging of network code.

In order to get QEMU networking to function, it is necessary to create
TAP interfaces, and, optionally, to bridge them to real networks.  A
development machine typically will have several TAP interfaces, and
more can be created.  Generally, each developer should have a TAP
interface of his or her own.  In the following, we will use our
development machine, newskysaw, as an example.

To set up a TAP interface on newskysaw, the following command is used:
\begin{verbatim}
/root/util/tap_create tapX
\end{verbatim}

When QEMU runs with a tap interface, it will use /etc/qemu-ifup to
bring up the interface.  On newskysaw, /etc/qemu-ifup looks like this:

\begin{verbatim}
#!/bin/bash
echo "Executing /etc/qemu-ifup - no external bridging"
echo "Bringing up $1 for bridged mode..."
NET=`echo $1 | cut -dp -f2` 
sudo /sbin/ifconfig $1 172.2${NET}.0.1 up
sleep 2
\end{verbatim}

The interface tap$N$ is brought up with the IP address 172.2$N$.0.1.
ifconfig will also create a routing rule that sends 172.2$N$.0.1/16
traffic to tap$N$.  The upshot is that if the code running in QEMU
uses an IP address in this network (for example: 172.2$N$.0.2), you
will be able to talk to it from newskysaw.  For example, from
newskysaw, if you ping 172.21.0.2, the packet (and ARP) will go out via
tap1.  The source address will appear to be 172.21.0.1.  The QEMU
machine will see these packets on its interface, and the software
controlling its interface can respond to 172.21.0.1.  

This form of networking is local to the machine.  You can also bridge
a TAP interface with a physical interface.  The result of this is that
a packet sent on it will be sent on the physical interface.  To do
this requires more effort (and is not set up by default on newskysaw).
As an example, consider that on newskysaw, the physical interface eth1
is connected to a private network switch to which the lab test
computers (v-test-amd, v-test-amd2, etc.) are connected.  To bridge,
for example, tap10, to this interface, you would do the following
(with root's help):
\begin{enumerate}
\item You need to bring up eth1 (ifconfig eth1 up {\em address}
netmask {\em mask}).  It is important that the address and mask you
choose are appropriate for the network eth1 is connected to.
\item You would bring up tap10 without an address:  /sbin/ifconfig
tap10 up
\item You would bridge tap10 and eth1:  /usr/sbin/brctl addif br0
tap10; /usr/sbin/brctl addif eth1.  This assumes that br0 was
previously created. 
\end{enumerate}

Bridging tap$N$ with eth1 will only work (where ``work'' means sending
a packet on the network and making the packet visible on localhost) if
the IP address in the code running in QEMU is set correctly.  This
means that it needs to be set to correspond to the network of eth1).  
For the newskysaw configuration, this is a 10-net address.


\subsection{Configuring Kitten}

Kitten needs to be explicitly configured to use networking. Currently
only a subset of the networking configurations are supported. To
enable an Ethernet network you should enable the following options:

\begin{itemize}
\item Enable TCP Support
\item Enable UDP Support
\item Enable socket API
\item Enable ARP support
\end{itemize}

The other options are not supported, and enabling them will probably
break the kernel compilation.

To allow Kitten to communicate with the QEMU network card you also
need to enable the appropriate device driver: \newline
\verb.NE2K Device Driver (rtl8139).

The driver then needs to be listed as a Kernel Command Line argument
in the {\em ISOIMAGE configuration}. To do this add
\verb.net=rtl819. to the end of the argument string.

Kitten currently does not support the dynamic assignment or IP
addresses at runtime. Because of this it is necessary to hardcode the
IP address into the device driver. For the rtl8139 network driver look
in the file {\em drivers/net/ne2k/rtl8139.c} for the function
\verb.rtl8139_init..

There should be a block of code that looks like the following:
\begin{verbatim}
  struct ip_addr ipaddr = { htonl(0 | 10 << 24 | 0 << 16 | 2 << 8 | 16 << 0) }; 
  struct ip_addr netmask = { htonl(0xffffff00) }; 
  struct ip_addr gw = { htonl(0 | 10 << 24 | 0 << 16 | 2 << 8 | 2 << 0) };
\end{verbatim}

This sets the IP address as 10.0.2.16, netmask 255.255.255.0 and
gateway address 10.0.2.2. Change these assignments to match your configuration.


\paragraph*{Kitten as the Guest OS}

When running Kitten as a VM, the above applies except that you will
want to enable the {\em VMNET} device driver instead of the {\em rtl8139}.


\subsection{Running with networking}

\paragraph*{TAP Interface}
Running with a TAP interface provides either local or global
connectivity (depending on how the TAP interface is configured and/or
bridged).  From the perspective of the QEMU command line, both look
the same, however.  You simply add something like this to the command
line:
\begin{verbatim}
-net tap,ifname=tap2 -net nic,model=rtl8139
\end{verbatim}
The first \verb.-net. option indicates that you want to use a tap
interface, specifically \verb.tap2..   The second \verb.-net. option
specifies that this interface will appear to code in the QEMU machine
to be a network interface card of the specific model RTL8139.  Note
that this is a model for which we have a driver.  If tap2 were
bridged, we'd get global connectivity.  If not, we would just get
local connectivity.  


\paragraph*{Redirection}
It is also possible to achieve limited local connectivity even if you
have no TAP support on your development machine.  In redirection, QEMU
essentially acts as a proxy, translating TCP or other connections and
low-level packet operations on the network interface in the QEMU
machine.  For example, the following options will redirect the host's
9555 port to the QEMU machine's 80 port:
\begin{verbatim}
-net user -net nic,model=rtl8139  -redir tcp:9555:10.10.10.33:80
\end{verbatim}
The first \verb.-net. option indicates that we are using user-level
networking (proxying).  The second \verb.-net. option indicates that
this user-level network will appear in the QEMU machine as an RTL8139
network card.   The \verb.-redir. option indicates that connections on
localhost:9555 will be translated into equivalent packet exchanges on
the RTL8139 card in the QEMU machine.  However, we have to tell QEMU
which IP address and port to use on the QEMU machine's side.  This is
what the 10.10.10.33 address, and port 80 are.  In the example, if you
access port 9555 on localhost, say with:
\begin{verbatim}
telnet localhost 9555
\end{verbatim}
The packets that appear in the QEMU machine will be bound for
10.10.10.33, port 80.  Within the QEMU machine, your RTL8139 interface
had better then be up on that address. 

QEMU has many options to build up virtual or real networking. See
http://www.h7.dion.ne.jp/$\sim$qemu-win/HowToNetwork-en.html for more
information.


For more questions, talk to Jack, Lei,  or Peter.

\end{document}
